\chapter{Mecanismos de Segurança}

O eBPF não é apenas uma ferramenta de observação passiva: integrado a subsistemas
de segurança consolidados do kernel Linux, ele viabiliza o controle ativo de
processos, redes e recursos do sistema com um nível de programabilidade que
abordagens tradicionais não oferecem \cite{linuxkernel2024}.
Este capítulo apresenta os principais mecanismos pelos quais o eBPF interage com
a infraestrutura de segurança do kernel.

% ----------------------------------------------------------------
\section{eBPF LSM}
% ----------------------------------------------------------------

O \textit{Linux Security Module} (LSM) é uma interface do kernel que define um
conjunto de \textit{hooks} em pontos de decisão de segurança — abertura de
arquivos, criação de processos, operações de rede, acesso a capacidades do sistema,
entre outros — e permite que módulos de segurança interceptem e influenciem essas
operações \cite{linuxkernel2024}.
Mecanismos tradicionais como SELinux e AppArmor são implementados como módulos
LSM, porém requerem configuração complexa e não podem ser adaptados em tempo de
execução.

A partir do kernel 5.7, o tipo de programa \texttt{BPF\_PROG\_TYPE\_LSM} permite
anexar programas eBPF diretamente aos \textit{hooks} LSM, criando módulos de
segurança totalmente programáveis \cite{linuxkernel2024}.
Ao retornar um valor de erro (por exemplo, \texttt{-EPERM}), o programa bloqueia
a operação antes de sua conclusão; ao retornar zero, permite que ela prossiga.
Essa decisão ocorre no próprio kernel, com latência de microssegundos, sem
qualquer envolvimento do espaço de usuário no caminho crítico
\cite{gregg2019}.

Entre os \textit{hooks} LSM mais relevantes para segurança estão:

\begin{itemize}
  \item \texttt{bpf\_lsm\_file\_open} — intercepta a abertura de arquivos,
        permitindo negar acesso com base em atributos do processo solicitante ou
        do arquivo;
  \item \texttt{bpf\_lsm\_bprm\_check\_security} — avalia a execução de um novo
        programa, viabilizando \textit{allowlists} de binários autorizados;
  \item \texttt{bpf\_lsm\_socket\_connect} — controla o estabelecimento de
        conexões de rede por processo, útil para isolar serviços comprometidos;
  \item \texttt{bpf\_lsm\_task\_kill} — regula o envio de sinais entre processos,
        prevenindo interferências maliciosas entre contêineres.
\end{itemize}

Uma vantagem significativa do LSM BPF em relação ao SELinux e ao AppArmor é a
possibilidade de combinar múltiplos programas no mesmo \textit{hook} — graças ao
mecanismo \textit{BPF link} — e de atualizar as políticas dinamicamente, sem
reiniciar serviços ou carregar novos módulos de kernel \cite{linuxkernel2024}.

% ----------------------------------------------------------------
\section{\textit{Cgroup hooks}}
% ----------------------------------------------------------------

\textit{Control groups} (cgroups) são o mecanismo do kernel Linux para organizar
processos em hierarquias e aplicar limites de recursos (CPU, memória, I/O de
rede e disco) a cada grupo \cite{linuxkernel2024}.
Em ambientes de contêineres, cada contêiner é tipicamente mapeado a um cgroup
distinto, tornando os \textit{hooks} de cgroup um ponto natural de aplicação de
políticas de segurança com escopo por contêiner \cite{gregg2020}.

O eBPF oferece um conjunto de tipos de programa associados a cgroups que permitem
interceptar operações antes que elas aconteçam:

\textbf{\texttt{BPF\_PROG\_TYPE\_CGROUP\_SKB}.}
Filtra pacotes na entrada ou saída de um cgroup específico.
Permite implementar políticas de rede por contêiner sem depender de regras de
\textit{iptables} ou \textit{nftables}, com menor overhead e maior expressividade
\cite{linuxkernel2024}.

\textbf{\texttt{BPF\_PROG\_TYPE\_CGROUP\_SOCK} e variantes.}
Controla a criação e o uso de sockets por processos do cgroup.
Pode bloquear a criação de sockets de determinados protocolos ou famílias de
endereço, restringindo efetivamente a superfície de ataque de rede de um
contêiner \cite{gregg2019}.

\textbf{\texttt{BPF\_PROG\_TYPE\_CGROUP\_DEVICE}.}
Regula o acesso a dispositivos (\texttt{/dev/*}) por processos do cgroup,
substituindo com vantagens a lista de dispositivos do cgroups v1
\cite{linuxkernel2024}.

\textbf{\texttt{BPF\_PROG\_TYPE\_CGROUP\_SYSCTL}.}
Intercepta leituras e escritas em parâmetros \texttt{sysctl} realizadas por
processos do cgroup, impedindo que um contêiner altere configurações globais do
kernel que afetariam outros contêineres \cite{linuxkernel2024}.

A granularidade por cgroup torna esses \textit{hooks} particularmente adequados
para ambientes Kubernetes, onde políticas de segurança precisam ser aplicadas
individualmente a cada \textit{pod} sem interromper os demais
\cite{cassagnes2020}.

% ----------------------------------------------------------------
\section{Seccomp e eBPF}
% ----------------------------------------------------------------

O \textit{secure computing mode} (seccomp) é um mecanismo do kernel Linux que
restringe o conjunto de chamadas de sistema que um processo pode invocar
\cite{linuxkernel2024}.
No modo \texttt{SECCOMP\_MODE\_FILTER}, introduzido no kernel 3.5, um programa BPF
clássico (cBPF) define a política de filtragem — especificando quais chamadas de
sistema são permitidas, bloqueadas ou encaminhadas para análise em espaço de usuário
\cite{mccanne1993}.

Com o eBPF, dois aspectos dessa integração evoluíram significativamente:

\textbf{Geração programática de filtros seccomp.}
Ferramentas modernas como \texttt{libseccomp} compilam políticas expressas em alto
nível para o formato cBPF exigido pelo kernel.
Pesquisas recentes exploram o uso de eBPF para gerar esses filtros de forma
dinâmica a partir de perfis de comportamento observados em execução — em vez de
listas estáticas definidas em tempo de desenvolvimento \cite{gregg2019}.

\textbf{Auditoria de violações seccomp com eBPF.}
Quando um filtro seccomp retorna \texttt{SECCOMP\_RET\_TRACE} ou
\texttt{SECCOMP\_RET\_LOG}, a chamada de sistema bloqueada pode ser capturada por
um programa eBPF conectado ao \textit{tracepoint} \texttt{syscalls:sys\_enter\_*},
permitindo registrar tentativas de evasão com contexto completo — processo, PID,
UID, argumentos — sem modificar o filtro seccomp original \cite{linuxkernel2024}.

A combinação de seccomp com LSM BPF representa a abordagem de defesa em
profundidade (\textit{defense in depth}) recomendada para ambientes de produção:
o seccomp restringe o conjunto de chamadas de sistema acessíveis, enquanto o LSM
BPF impõe controles adicionais sobre o que pode ser feito com as chamadas
permitidas \cite{beyer2016}.
Juntos, os dois mecanismos reduzem drasticamente a superfície de ataque exposta
por um processo comprometido \cite{gregg2020}.

% ----------------------------------------------------------------
\section{Controle de processos e acesso a recursos}
% ----------------------------------------------------------------

Além dos mecanismos anteriores — orientados a interfaces específicas como sockets,
arquivos e chamadas de sistema —, o eBPF oferece visibilidade e controle sobre o
ciclo de vida de processos e o uso de recursos do sistema operacional de forma
holística \cite{gregg2019}.

\textbf{Rastreamento de ancestralidade de processos.}
Utilizando mapas BPF do tipo \texttt{BPF\_MAP\_TYPE\_HASH}, é possível manter em
tempo real uma tabela de processos ativos, registrando PID, PPID, UID, executável
e cgroup de cada processo \cite{linuxkernel2024}.
Essa árvore de processos viabiliza a detecção de \textit{process injection} — em
que um atacante insere código em um processo legítimo —, pois alterações inesperadas
na identidade ou no comportamento de um processo tornam-se visíveis
\cite{majors2022}.

\textbf{Monitoramento de uso de capacidades (\textit{capabilities}).}
O modelo de capacidades do Linux substitui o binômio root/não-root por um conjunto
granular de privilégios (\texttt{CAP\_NET\_ADMIN}, \texttt{CAP\_SYS\_PTRACE},
\texttt{CAP\_DAC\_OVERRIDE}, entre outras) \cite{linuxkernel2024}.
Programas eBPF conectados ao \textit{hook} LSM \texttt{bpf\_lsm\_capable} podem
registrar toda tentativa de uso de capacidade, identificando processos que operam
com mais privilégios do que o esperado — um indicador comum de escalonamento de
privilégios \cite{gregg2020}.

\textbf{Controle de acesso a namespaces.}
Namespaces do Linux isolam recursos como PID, rede, sistema de arquivos e
identificadores de usuário entre grupos de processos \cite{linuxkernel2024}.
Tentativas de um processo sair do seu namespace — por exemplo, via
\texttt{setns(2)} ou \texttt{unshare(2)} — podem ser monitoradas e bloqueadas
por programas eBPF, prevenindo fugas de contêineres (\textit{container escapes})
\cite{cassagnes2020}.

\textbf{Observabilidade de acesso a arquivos e descritores.}
O \textit{tracepoint} \texttt{sched\_process\_exit}, combinado com um mapa de
estado de arquivos abertos, permite detectar processos que terminam sem fechar
descritores sensíveis — padrão observado em algumas classes de \textit{malware}
que abusam de herança de descritores entre processos \cite{gregg2019}.

\textbf{Limitação dinâmica de recursos.}
Em conjunto com os \textit{cgroup hooks} descritos na seção anterior, programas
eBPF podem implementar cotas dinâmicas de chamadas de sistema por processo — por
exemplo, limitando a taxa de chamadas \texttt{fork(2)} para mitigar ataques de
exaustão de processos (\textit{fork bombs}) — sem necessidade de reconfigurar
os limites de cgroup em espaço de usuário \cite{linuxkernel2024}.

A combinação desses controles com os mecanismos de observabilidade descritos no
capítulo anterior cria uma camada de segurança coesa: eventos suspeitos são
detectados, correlacionados e, quando necessário, bloqueados em uma única
infraestrutura baseada em eBPF, sem a sobrecarga de agentes externos e com
latência compatível com ambientes de produção \cite{majors2022}.